\documentclass[10pt,a4paper]{amsart}
\usepackage[margin=19mm]{geometry}
\usepackage{amsmath,amssymb,mathtools}
\usepackage[scr=rsfs,cal=euler]{mathalfa}
\usepackage{xfrac}
\usepackage[unicode,hidelinks,bookmarks=false]{hyperref}
\newcommand{\ie}{i.e.\ }
\newcommand{\cf}{cf.\ }
\newcommand{\resp}{respectively\ }
\newcommand{\Subsection}{\subsection}
\providecommand{\nobreakdash}{\nobreak\hbox{-}}
\setlength{\emergencystretch}{2em}
\multlinegap0pt
\usepackage{bm}
\usepackage{amssymb}
\usepackage{xcolor}
\usepackage{mathtools}
\usepackage{algorithm}\usepackage{algpseudocode}
\newtheorem{lemma}{Lemma}
\title{MAGPIE: Multilevel--Adaptive--Guided solver for ptychographic phase retrieval}
\author{Borong Zhang, Qin Li, Zichao Wendy Di}
\date{Source: arXiv:2504.10118v7 (CC BY 4.0)}
\begin{document}
\maketitle

\noindent\emph{Source and licence.} MAGPIE: Multilevel--Adaptive--Guided solver for ptychographic phase retrieval. Version arXiv:2504.10118v7 (CC BY 4.0), distributed by arXiv under the Creative Commons Attribution 4.0 International licence, http://creativecommons.org/licenses/by/4.0/, as recorded in the arXiv licence metadata for this exact version and shown on the abstract page https://arxiv.org/abs/2504.10118v7. Full licence notice: \texttt{licenses/arxiv-2504.10118-CC-BY-4.0.txt}.\par\smallskip

\noindent\textit{Editorial scope.} Counted unit: the article's lemma lem:commutative (source0.tex lines 1045--1081). The article's complete original proof of it is delivered with it (source0.tex lines 1162--1197, one \texttt{proof} environment of 36 lines, which the article places away from the statement). No mathematics is written, repaired or re-derived: the statement and the proof are the article's own lines, in the article's own notation, and no retained line is substituted or re-flowed.  Retained uncounted as the source-local context the proof needs: lines 1083--1092; lines 1094--1103.  Nothing was omitted from any retained span, so the package carries no omission record.  No retained line contains a citation, so the wrapper carries no bibliography and no bibliography entry.  No retained line contains a figure environment, an image-inclusion command or drawing code, so no image is delivered and the package declares no asset dependency.  Editorial formatting only, in addition: an AMS \texttt{amsart} preamble that restates the article's own declarations and macros used by the retained lines, namely the environments \texttt{lemma} on separate counters, together with the standard packages those lines need.  The retained environments print in this document as Lemma 1 in order of appearance, whereas the article prints the same items with the numbering of its own full document.  The complete native source of the article as distributed in the arXiv e-print for this exact version is bundled unchanged as \texttt{sources/176498/source0.tex}.
\begin{lemma}\label{lem:unary}
    For any function $f:\mathbb{C}\to\mathbb{C}$, we define the operators
\begin{equation*}
    \bm{f}_H:\mathbb{C}^{n_H^2}\to\mathbb{C}^{n_H^2}\quad\text{and}\quad \bm{f}_h:\mathbb{C}^{n_h^2}\to\mathbb{C}^{n_h^2}\,,
\end{equation*}
to be the element-wise applications of $f$, i.e., for any input vector, the operator applies $f$ to every entry. Then, for any $\bm{B}\in\mathbb{C}^{n_H^2}$,
\begin{equation*}
    \bm{f}_h\left(\bm{I}_H^h\bm{B}\right)= \bm{I}_H^h\bm{f}_H(\bm{B})\,.
\end{equation*}
\end{lemma}

\begin{lemma}\label{lem:binary}
    For any $\bm{B}$, and $\bm{C}\in\mathbb{C}^{n_H^2}$ and any binary function $g:\mathbb{C}\otimes\mathbb{C}\to\mathbb{C}$, we define the operators
\begin{equation*}
    \bm{g}_H:\mathbb{C}^{n_H^2}\otimes\mathbb{C}^{n_H^2}\to\mathbb{C}^{n_H^2}\quad\text{and}\quad \bm{g}_h:\mathbb{C}^{n_h^2}\otimes\mathbb{C}^{n_h^2}\to\mathbb{C}^{n_h^2}\,,
\end{equation*}
to be the element-wise applications of $g$. Then, for any $\bm{B},\bm{C}\in\mathbb{C}^{n_H^2}$,
\begin{equation*}
    \bm{g}_h\left(\bm{I}_H^h\bm{B},\bm{I}_H^h\bm{C}\right)= \bm{I}_H^h\bm{g}_H(\bm{B},\bm{C})\,.
\end{equation*}
\end{lemma}

\begin{lemma}\label{lem:commutative}
For any $\bm{A}\in\mathbb{C}^{n_h^2}$ and $\bm{B}\in\mathbb{C}^{n_H^2}$, we have
    \begin{equation*}
        \bm{B}\odot \bm{I}_h^H\bm{A} = \bm{I}_h^H((\bm{I}_H^h\bm{B})\odot\bm{A})\,.
    \end{equation*}
Note that when $\bm{A}=\bm{J}$, a vector of ones, we have 
\begin{equation*}
    \bm{B} = \bm{I}_h^H\bm{I}_H^h\bm{B}\implies \bm{I}_h^H\bm{I}_H^h = \bm{I},\text{ the identity matrix}\,.
\end{equation*}

    \begin{proof}
    Note that using the binning strategy, each column of $\bm{I}_h^H$ is non-zero only at one entry, i.e.,
  \begin{equation*}
    (\bm{I}_h^H)_{ij}(\bm{I}_h^H)_{kj} = \frac{1}{4}(\bm{I}_h^H)_{ij}\delta_{ik}\,,
  \end{equation*}
  and by definition, we have
  \begin{equation*}
      \bm{I}_H^h = 4(\bm{I}_h^H)^\top\,.
  \end{equation*}
    For all entries in the $\texttt{LHS}$ and $\texttt{RHS}$, we have 
    \begin{equation*}
            \texttt{LHS}_i = \bm{B}_i(\bm{I}_h^H\bm{A})_i = \bm{B}_i\sum_{j}(\bm{I}_h^H)_{ij}\bm{A}_j\,.
    \end{equation*}
    \item 
    \begin{equation*}
        \begin{aligned}
            \texttt{RHS}_i &= (\bm{I}_h^H[(\bm{I}_H^h\bm{B})\odot\bm{A}])_i\,,\\
            &= \sum_{j}(\bm{I}_h^H)_{ij}[(\bm{I}_H^h\bm{B})\odot\bm{A}]_j\,,\\
            &= \sum_{j}(\bm{I}_h^H)_{ij}\left[\sum_{k}(\bm{I}_H^h)_{jk}\bm{B}_k\right]\bm{A}_j\,,\\
            &= \sum_{k}\bm{B}_k\sum_{j}(\bm{I}_h^H)_{ij}(\bm{I}_H^h)_{jk}\bm{A}_j\,,\\
            &= \sum_{k}\bm{B}_k\sum_{j}(\bm{I}_h^H)_{ij}4(\bm{I}_h^H)_{kj}\bm{A}_j\,,\\
             &= \sum_{k}\bm{B}_k\sum_{j}(\bm{I}_h^H)_{ij}\delta_{ik}\bm{A}_j\,,\\
             &= \bm{B}_i\sum_{j}(\bm{I}_h^H)_{ij}\bm{A}_j\,.
        \end{aligned}
    \end{equation*}
    \end{proof}
\end{lemma}

\begin{proof} We provide proofs for the statements in the same order. 
\begin{itemize}
    \item By Lemma~\ref{lem:commutative} and Lemma~\ref{lem:unary}, we have
    \begin{equation*}
        \begin{aligned}
            \bm{I}_h^H{\bm{W}_{\bm{z}}} &= \bm{I}_h^H\left( \frac{|\bm{Q}|^2}{\bm{I}_H^h\bm{I}_h^H|\bm{Q}|^2}\right)\,,\\
            &= \bm{I}_h^H\left( \bm{I}_H^h\left(\frac{1}{\bm{I}_h^H|\bm{Q}|^2}\right)\odot|\bm{Q}|^2\right)\,,\\
            &= \left(\frac{1}{\bm{I}_h^H|\bm{Q}|^2}\right)\odot\bm{I}_h^H|\bm{Q}|^2\,,\\
            &= \bm{J}\,,\quad\text{a vector of ones}\,.
        \end{aligned}
    \end{equation*}
    For the sake of a contradiction, we assume that $(\bm{W}_{\bm{z}})_r > 4$ at some index $r$. Since $\bm{W}_{\bm{z}}\geq 0$, we have $\bm{I}_h^H{\bm{W}_{\bm{z}}}\neq \bm{J}$ because in the bin of index $r$, the average is greater than 1. That is a contradiction to the first statement. Therefore, $\bm{W}_{\bm{z}}$ element-wise bounded by $4$.
    \item First, we find that
    \begin{equation*}
        \left|\frac{(\bm{I}_H^h\bm{Q}_H)\odot\bm{W}_{\bm{z}}}{\bm{Q}}\right| =  \frac{|(\bm{I}_H^h\bm{Q}_H)|\odot|\bm{Q}|^2}{|\bm{Q}|\odot\left(\bm{I}_H^h\bm{I}_h^H|\bm{Q}|^2\right)} = \frac{\left(\bm{I}_H^h|\bm{Q}_H|\right)\odot|\bm{Q}|}{\bm{I}_H^h\bm{I}_h^H|\bm{Q}|^2}\leq \frac{\left(\bm{I}_H^h\bm{I}_h^H|\bm{Q}|\right)\odot|\bm{Q}|}{\bm{I}_H^h\bm{I}_h^H|\bm{Q}|^2} \,.
    \end{equation*}
    Then, by Lemma~\ref{lem:commutative}, Lemma~\ref{lem:unary}, and Lemma~\ref{lem:binary}, we have
    \begin{equation*}
        \begin{aligned}
            \bm{I}_h^H\left|\frac{(\bm{I}_H^h\bm{Q}_H)\odot\bm{W}_{\bm{z}}}{\bm{Q}}\right|&\leq\bm{I}_h^H\left(\frac{\left(\bm{I}_H^h\bm{I}_h^H|\bm{Q}|\right)\odot|\bm{Q}|}{\bm{I}_H^h\bm{I}_h^H|\bm{Q}|^2}\right)\,,\\&=\bm{I}_h^H\left(\bm{I}_H^h\left(\frac{\bm{I}_h^H|\bm{Q}|}{\bm{I}_h^H|\bm{Q}|^2}\right)\odot|\bm{Q}|\right)\,,\\
            &=\left(\frac{\bm{I}_h^H|\bm{Q}|}{\bm{I}_h^H|\bm{Q}|^2}\right)\odot\bm{I}_h^H|\bm{Q}|\,,\\
            &=\frac{(\bm{I}_h^H|\bm{Q}|)^2}{\bm{I}_h^H|\bm{Q}|^2}\leq 1\,,
        \end{aligned}
    \end{equation*}
    where the last element-wise inequality follows from Jensen's inequality applied to each bin. Hence, by exactly the same argument as the proof of the previous statement, we have
    \begin{equation*}
        \bm{I}_h^H\left|\frac{(\bm{I}_H^h\bm{Q}_H)\odot \bm{W}_{\bm{z}}}{\bm{Q}}\right|\leq 1 \implies \left|\frac{(\bm{I}_H^h\bm{Q}_H)\odot\bm{W}_{\bm{z}}}{\bm{Q}}\right|\leq 4\,,
    \end{equation*}
    where the inequalities are satisfied element-wise. We note the sharp bound for the last inequality is $3/2$, but its more involved derivation is omitted here.
    \item By applying Jensen's inequality to each bin, 
    \begin{equation*}
        |\bm{Q}_H|^2 \leq\left(\bm{I}_h^H|\bm{Q}|\right)^2\leq \bm{I}_h^H |\bm{Q}|^2\implies \frac{|\bm{Q}_H|^2}{\bm{I}_h^H|\bm{Q}|^2}
        \leq 1\,.
    \end{equation*}
\end{itemize}
\end{proof}

\end{document}
